% !TEX root = userguide.tex

\def\headdate{23rd June 2024}

\section{Elastoviscoplastic flow using the Saramito models}
\label{sec:evp_flow}

This section describes some examples of elastoviscoplastic (evp) flow of 
Saramito constitutive models as described in
Appendix~\ref{chap:viscoelasticmodels}.
 
 \subsection{Start-up of a Poiseuille flow of an elastoviscoplastic material}
 
The example file is \texttt{poiseuille2.f90} and can
be obtained by typing at the
command line:
\begin{quote}
   \texttt{getexample poiseuille}
\end{quote}
The example uses a unit square domain with two elements in $x$-direction and
periodical boundary conditions between in and outflow.
The flow is generated by imposing a constant flow rate using a constraint.
The time integration is second-order implicit Gear with conformation prediction
(Eq.~\eqref{eq:Gearimp}).
The model is a single-mode Oldroyd-B/UCM, Giesekus, PTT-linear or
PTT-exponential model with adapted $\lambda$ to create an elastoviscoplastic model, according to Saramito (the SRM1 and SRM2 models \cite{saramito_new_2007,saramito_new_2009}). See also Sec.~\ref{sec:alammodels}.
Optionally, transient inertia terms (see Sec.~\ref{sec:navierstokes})
can be included. For the time-discretization of the transient inertia terms we
use the BDF2 scheme (See~Eq.\eqref{eq:NS_Kar2}).
You can run the example by
\begin{quote}
   \texttt{poiseuille2 < input\_poiseuille2.txt}\\
   \texttt{poiseuille\_post}
\end{quote}
Inspect the source files to see which data and plotfiles are generated.

The example file \texttt{poiseuille3.f90} contains the same problem as \texttt{poiseuille2.f90}, but now using the CDT formulation (see Sec~\ref{sec:CDTformulation}). You can run the example by
\begin{quote}
   \texttt{poiseuille3 < input\_poiseuille3.txt}\\
   \texttt{poiseuille\_post}
\end{quote}
Inspect the source files to see which data and plotfiles are generated.

In Figure~\ref{fig:poiseuille_evp} results are shown of a steady state profile across the half-height of the channel.
 Two regions can be seen: an elastic solid region in the middle of the channel that moves with a constant velocity and wall region, where the material yields and becomes a viscoelastic fluid. Some notes:
\begin{itemize}
 \item A nice linear continuous shear stress profile is present across the two regions, indicating a steady state profile.
 \item In the solid elastic region the shear rate is zero, but the stress, and thus the strain, is not.
 \item The von Mises equivalent shear stress $\tau_\text{d}$ consists of both a shear and normal stress part. Therefore, yielding starts at a $y$ value where $\tau_{xy}<\tau_\text{y}$.
 \item At the $y$ value where yielding starts, there is a jump in shear rate. Although the shear stress $\tau_{xy}$ is continuous, the normal stress $\tau_{xx}$ is not. As a result, there is a jump in the von Mises stress also.
\end{itemize}
\begin{figure}[htbp!]
\begin{center}
    \includegraphics[scale=1.0]{poiseuille_evp}
\end{center}
   \caption{Velocity and von Mises equivalent shear stress profile in a steady Poiseuille flow of a SRM2 evp model using the example \texttt{poiseuille2.f90}.
   The center of the channel is at $y=0$ and the wall at $y=1$. The imposed average velocity $U_\text{avg}=1$. Model parameters: neo-Hookean solid below the yield point and a UCM based viscoelastic model beyond the yield point, $G=10$, $\tau_y=10$, $K=1$, $n=0.5$. Number of elements across the channel half-width is 200.}
    \label{fig:poiseuille_evp}
\end{figure}

\subsection{Axi-symmetric problem: flow around a sphere in
 a cylindrical tube filled with an elastoviscoplastic fluid}
\label{sec:sphereimplicitevp}

The file \texttt{sphere27.f90} in the example `sphere' is similar to
\texttt{sphere22.f90} (see Sec.~\ref{sec:sphereimplicit}, but now using an elastoviscoplastic model
(Sec.~\ref{sec:alammodels}). The input file is \texttt{input\_sphere27.txt}.
The example file \texttt{sphere28.f90} contains the same problem as
\texttt{sphere27.f90}, but now using the CDT formulation (see
Sec~\ref{sec:CDTformulation}). The input file is \texttt{input\_sphere28.txt}.
	
The flow becomes unsteady with a clear periodicity beyond a certain yield stress value. This needs further investigation.

 \subsection{Start-up of a Poiseuille flow of an elastoviscoplastic material with Drucker-Prager plasticity}
 
The example file is \texttt{poiseuille4.f90} and can be obtained by typing at the command line:
\begin{quote}
   \texttt{getexample poiseuille}
\end{quote}
The example \texttt{poiseuille4.f90} is similar to \texttt{poiseuille2.f90}, but now the model is replaced by the Saramito model based on Drucker-Prager plasticity \cite{saramito_new_2021} (the SRM3 model).
You can run the example by
\begin{quote}
   \texttt{poiseuille4 < input\_poiseuille4.txt}\\
   \texttt{poiseuille\_post}
\end{quote}
Inspect the source files to see which data and plotfiles are generated.

In Figure~\ref{fig:poiseuille_evp2} results are shown of profiles across the half-height of the channel at time $t=200$.
The flow is near a steady state.
Three regions can be seen: an elastic solid region in the middle of the channel (regime 1, ``sticking'') that moves with a constant velocity, a second region (regime 2, ``sliding'') and a wall region (regime 3, ``losing contact'').
\begin{itemize}
\item Regime 1 is elastic and the stress is a pure elastic stress such that it is below the yield surface.
\item Regime 2 is governed by the change in $\tr\ten\tau$ and the pressure $p$ adjusts accordingly to obtain a smooth Cauchy stress tensor.
\item Regime 2 itself seems to be splitted into two parts: a part where the material remains a solid and a part where the yielding takes place. The first part is actually very close to the yield surface and can be considered to be solid. Extending the simulation time, shows that more and more points in this part move over to regime 1. Also the transition from solid to yielding becomes even more sharp. The yielding starts, at the position where the shear stress $\tau$ reaches the value \cite{Hulsennote500}:
\[
   \tau_0 = \frac12(\mu G + \tau_\text{y})\dfrac{\sqrt{1-\frac43\mu^2}}{1-\frac13\mu^2}
\]
\end{itemize}

\begin{figure}[htbp!]
\begin{center}
    \includegraphics[scale=0.75]{poiseuille_evp2}
\end{center}
   \caption{Velocity, vorticity, stress and regime profiles in a steady Poiseuille flow of a SRM3 evp model using the example \texttt{poiseuille4.f90}. Time $t=100$.
The center of the channel is at $y=0$ and the wall at $y=1$. The imposed average velocity $U_\text{avg}=1$. Model parameters: neo-Hookean solid below the yield point with Drucker-Prager yield surface. $G=10$, $\lambda=0.1$, $\tau_\text{y}=0.2$ and $\mu=0.2$. Number of elements across the channel half-width is 200. Note: the regime (1, 2 or 3) has been multiplied with 15 to better see the different regimes.}
    \label{fig:poiseuille_evp2}
\end{figure}



 




